\documentclass[12pt, a4paper]{article}
\usepackage[T1]{fontenc}
\usepackage{lmodern}
\usepackage{amsmath, amsthm, amssymb, mathtools}
\usepackage[margin=1in]{geometry}
\usepackage{xr-hyper}
\usepackage[colorlinks=true,linkcolor=blue,citecolor=blue]{hyperref}
\usepackage{cleveref}
\usepackage{graphicx, subcaption}
\usepackage{natbib}
\usepackage{booktabs}
\usepackage{enumerate}

% Code references: scaled monospace with line-breaking
\newcommand{\code}[1]{\texttt{\small #1}}

\newtheorem{theorem}{Theorem}[section]
\newtheorem{lemma}[theorem]{Lemma}
\newtheorem{proposition}[theorem]{Proposition}
\newtheorem{corollary}[theorem]{Corollary}
\newtheorem{definition}{Definition}[section]
\newtheorem{remark}{Remark}[section]
\newtheorem{conjecture}[theorem]{Conjecture}
\newtheorem{derivation}[theorem]{Derivation}
\newtheorem{observation}[theorem]{Observation}
\numberwithin{equation}{section}

\newcommand{\dd}{\mathrm{d}}
\newcommand{\seriestitle}{Why Rotating Fluids Make Polygons}


\externaldocument[I-]{../paper-1-mathematics/main}
\externaldocument[II-]{../paper-2-physics/main}
\externaldocument[III-]{../paper-3-gravity/main}
\externaldocument[V-]{../paper-4-planets/main}

\title{\seriestitle\\[6pt]\large Paper IV: The Havelock Field Theory}
\author{Gordon Speagle}
\date{}

\begin{document}
\maketitle

\begin{abstract}
We construct a complete quantum field theory on the Seifert manifold
$\mathbb{R} \times (\mathbf{H}^2 \times_N S^1)$ with Thurston
$\widetilde{\mathrm{SL}(2,\mathbb{R})}$ geometry, parameterised by a
single integer~$N$.
The theory contains gravity, a $\mathrm{U}(1)$ gauge field,
a charged scalar, and Dirac fermions, with all couplings determined
by the central charge $c = N^2$:
$G = 3/(2c)$, $\alpha = 6/(\pi c)$,
$\alpha/(8\pi G) = 1/(2\pi^2)$ (exact, $N$-independent).

The cosmological constant
$\Lambda = (N^2 - 16)/16$ is positive for $N \ge 5$,
with $\Lambda = 0$ at $N = 4$ (the gauge threshold where $j = 1$).
The Standard Model gauge group
$\mathrm{SU}(3) \times \mathrm{SU}(2) \times \mathrm{U}(1)$ emerges
from the polygon hierarchy:
$\mathrm{SU}(3)$ from the Galois cover at $N = 7$
($\operatorname{Gal}(K_7/\mathbb{Q}) = \mathbb{Z}/3\mathbb{Z}
= Z(\mathrm{SU}(3))$);
$\mathrm{SU}(2)$ from the Chern--Simons adjoint Wilson line at $N = 4$
($j = 1$, the adjoint of $\mathrm{SL}(2)$);
$\mathrm{U}(1)$ from Kaluza--Klein reduction on $S^1$.
The Weinberg angle $\sin^2\theta_W = 3/11 \approx 0.273$
(18\% from the experimental value~$0.231$, versus the SU$(5)$ GUT
prediction~$3/8 = 0.375$, which is 62\% away).

The Weyl anomaly $\delta_m = 0$ exactly on all constant-curvature
surfaces; the three-layer decomposition is exact with zero anomaly.
The theory is UV-complete through the orbifold CFT at $c = N^2$
(non-perturbative) and super-renormalisable after Kaluza--Klein
reduction to two dimensions.
\end{abstract}

%% ====================================================================
\section{Introduction: one integer, one geometry}
\label{sec:intro}

The preceding papers established that the Havelock eigenvalue
$\lambda_m = C_1(\xi) - m(N{-}m)/2$ governs both point vortex
polygon stability (Paper~II) and gravitational polygon stability
in $2{+}1$D (Paper~III), with the palindromic thresholds
encoding algebraic number fields (Paper~I).
Paper~III derived the vacuum Einstein equations from polygon stability
and identified the graviton at $N = 7$ ($j = 2$, the unique polygon
number with integer spin~$\ge 2$ at the critical mode).

This paper constructs the full quantum field theory.
The single input is an integer $N \ge 3$ (the polygon number $=$
the flux quantum $=$ the superselection sector).
Everything else follows:
\begin{enumerate}
\item \S\ref{sec:spacetime}: the spacetime
  $\mathbb{R} \times (\mathbf{H}^2 \times_N S^1)$
  with Thurston $\widetilde{\mathrm{SL}(2,\mathbb{R})}$ geometry;
\item \S\ref{sec:fields}: gravity, gauge, scalar, fermion, radion;
\item \S\ref{sec:couplings}: $G$, $\alpha$, $g_4$---all from $c = N^2$;
\item \S\ref{sec:lambda}: the cosmological constant
  $\Lambda = (N^2 - 16)/16$;
\item \S\ref{sec:anomaly}: the Weyl anomaly vanishes
  ($\delta_m = 0$ on constant-curvature backgrounds);
\item \S\ref{sec:gauge-group}: the Standard Model gauge group
  $\mathrm{SU}(3) \times \mathrm{SU}(2) \times \mathrm{U}(1)$;
\item \S\ref{sec:weinberg}: the Weinberg angle $\sin^2\theta_W = 3/11$;
\item \S\ref{sec:two-forces}: the two-force unification
  (gravity/colour $+$ electroweak);
\item \S\ref{sec:uv}: UV completion via the orbifold CFT.
\end{enumerate}

%% ====================================================================
\section{The spacetime: $\mathbb{R} \times (\mathbf{H}^2 \times_N S^1)$}
\label{sec:spacetime}

The four-dimensional spacetime is the product of a time axis with
a Seifert-fibred three-manifold:
\begin{equation}\label{eq:metric-4d}
  ds^2 = -dt^2 + ds^2_{\mathbf{H}^2}
  + (d\varphi + A_i\,dx^i)^2,
\end{equation}
where $ds^2_{\mathbf{H}^2} = d\rho^2 + \sinh^2\!\rho\;d\theta^2$
($K = -1$) and $A$ is a $\mathrm{U}(1)$ connection on the $S^1$ fiber
with flux
\begin{equation}\label{eq:flux}
  \frac{1}{2\pi}\oint dA = \frac{N}{2}.
\end{equation}
This is a Thurston geometry of type
$\widetilde{\mathrm{SL}(2,\mathbb{R})}$ \citep{Thurston1982}.
The Kaluza--Klein reduction on $S^1$ decomposes any field into
$N$ modes with masses
\begin{equation}\label{eq:kk-masses}
  \mu_m = |m - N/2|, \qquad m = 0, 1, \ldots, N{-}1.
\end{equation}
The critical mode $m = N/2$ (even $N$) is massless: it IS the
Havelock critical mode, and the Casimir $f(m,N) = m(N{-}m)/2$
IS the kinetic energy on $S^1$.

%% ====================================================================
\section{The field content}
\label{sec:fields}

The metric~\eqref{eq:metric-4d} yields, after KK reduction:
\begin{enumerate}
\item \textbf{Gravity.}
  The $2{+}1$D metric $g_{ij}$ on $\mathbf{H}^2$ is a
  Chern--Simons gauge field at level
  $k = c/12 = N^2/12$
  (Proposition~\ref{III-prop:graviton}, Paper~III).
\item \textbf{$\mathrm{U}(1)$ gauge field.}
  The off-diagonal component $A_\mu$ from the KK ansatz;
  coupling $\alpha = 6/(\pi c)$.
\item \textbf{Charged scalar.}
  $N$ KK modes $\Phi_m$ with masses $\mu_m = |m - N/2|$.
  The critical mode ($m = N/2$) is massless.
\item \textbf{Dirac fermion.}
  $N$ KK modes $\Psi_m$ with masses $\mu^f_m = |m - (N{-}1)/2|$
  (shifted by $1/2$ from the spin connection).
  Even~$N$: bosons have a massless mode, fermions do not
  $\implies$ SUSY broken by flux quantisation.
\item \textbf{Radion.}
  The $S^1$ modulus $\sigma$ (fiber radius), stabilised by
  the flux energy at $\sigma = 1$.
\end{enumerate}

%% ====================================================================
\section{The couplings: all from $c = N^2$}
\label{sec:couplings}

\begin{proposition}[Exact coupling relations]
\label{prop:couplings}
All couplings are determined by $c = 12\,b(N) = N^2 + O(N)$:
\begin{align}
  G &= \frac{3}{2c}, \label{eq:newton}\\
  \alpha &= \frac{6}{\pi c}, \label{eq:alpha}\\
  \frac{\alpha}{8\pi G} &= \frac{1}{2\pi^2}
    \qquad\text{(exact, $N$-independent)}. \label{eq:coupling-lock}
\end{align}
The quartic vertex from the Magri hierarchy:
$g_4 = 2N f^2/\pi$ at the critical mode.
\end{proposition}

\begin{proof}
$G$: from Brown--Henneaux, $c = 3\ell/(2G)$ with $\ell = 1$.
$\alpha$: from the KK gauge coupling on the unit-radius fiber.
The ratio~\eqref{eq:coupling-lock}: $\alpha/(8\pi G)
= [6/(\pi c)] / [8\pi \cdot 3/(2c)] = 6/(12\pi^2) = 1/(2\pi^2)$.
The $N$-dependence cancels algebraically.
\end{proof}

%% ====================================================================
\section{The cosmological constant}
\label{sec:lambda}

\begin{proposition}[Cosmological constant from the fiber]
\label{prop:lambda}
The effective two-dimensional cosmological constant is
\begin{equation}\label{eq:lambda}
  \Lambda = \frac{N^2 - 16}{16}.
\end{equation}
$\Lambda = 0$ at $N = 4$ (the gauge threshold, $j = 1$);
$\Lambda < 0$ for $N = 3$ (AdS);
$\Lambda > 0$ for $N \ge 5$ (de~Sitter).
\end{proposition}

\begin{proof}
The fiber curvature from the magnetic flux:
$R_{\varphi\varphi} = (N/2)^2/4 = N^2/16$.
The base curvature: $R_{\mathbf{H}^2} = -1$.
The effective cosmological constant:
$\Lambda = R_{\text{fiber}} - |R_{\text{base}}| = N^2/16 - 1
= (N^2 - 16)/16$.
\end{proof}

\noindent
The two thresholds have a unified interpretation:
\begin{center}
\small
\begin{tabular}{lcll}
\toprule
$N$ & $\Lambda$ & $j_{\mathrm{crit}}$ & Physical meaning \\
\midrule
$4$ & $0$ & $j = 1$ & Gauge boson; $\Lambda = 0$ boundary \\
$7$ & $33/16$ & $j = 2$ & Graviton; stability boundary $N_{\mathrm{crit}} = 7$ \\
\bottomrule
\end{tabular}
\end{center}

%% ====================================================================
\section{The Weyl anomaly vanishes}
\label{sec:anomaly}

\begin{theorem}[Zero anomaly on Einstein backgrounds]
\label{thm:zero-anomaly}
On any surface of constant curvature
(the plane, $\mathbf{S}^2$, $\mathbf{H}^2$, or any quotient
$\Gamma\backslash\mathbf{H}^2$), the constrained Havelock eigenvalue
satisfies
\[
  \lambda_m = C_1(\mathrm{surface}, R) - \frac{m(N{-}m)}{2}
\]
with $\delta_m = 0$ \emph{identically}.
The three-layer decomposition is exact with zero Weyl anomaly.
\end{theorem}

\begin{proof}
The Havelock identity (Paper~I, Theorem~\ref{I-thm:riemannian-havelock})
proves that the \emph{constrained} eigenvalue (on $\{J = \mathrm{const}\}$)
has a mode-independent confining potential~$C_1$.
The angular Hessian involves the csc$^2$ kernel from the logarithmic
singularity of the Green's function; the Havelock sum
(Paper~I, Lemma~\ref{I-lem:havelock}) converts this to $m(N{-}m)/2$.
No residual $\delta_m$ survives.
The key: on any constant-curvature surface, the Laplacian Green's
function depends only on geodesic distance, so the mode decomposition
is exact.
Quotients $\Gamma\backslash\mathbf{H}^2$ add a mode-\emph{independent}
correction $\delta C_1$ (from the Selberg trace formula,
Paper~I, Theorem~\ref{I-thm:index-pairing}), which shifts~$C_1$
but does not introduce mode-dependent anomalies.
\end{proof}

%% ====================================================================
\section{The Standard Model gauge group}
\label{sec:gauge-group}

The polygon hierarchy realises
$\mathrm{SU}(3) \times \mathrm{SU}(2) \times \mathrm{U}(1)$
through three independent mechanisms.

\subsection{$\mathrm{SU}(3)$ from the Galois cover at $N = 7$}

The trace field $K_7 = \mathbb{Q}(\cos 2\pi/7)$ has degree
$[K_7:\mathbb{Q}] = 3$.
The Galois group
$\operatorname{Gal}(K_7/\mathbb{Q}) \cong \mathbb{Z}/3\mathbb{Z}$
acts on the three palindromic mode pairs of the $N = 7$ polygon:
\begin{center}
\small
\begin{tabular}{cccc}
\toprule
Pair & Modes & Casimir $f$ & Galois orbit \\
\midrule
1 & $(m{=}1, m{=}6)$ & $3$ & $1 \to 2 \to 3 \to 1$ \\
2 & $(m{=}2, m{=}5)$ & $5$ & $2 \to 3 \to 1 \to 2$ \\
3 & $(m{=}3, m{=}4)$ & $6$ & $3 \to 1 \to 2 \to 3$ \\
\bottomrule
\end{tabular}
\end{center}
\noindent
The cyclic permutation identifies
$\mathbb{Z}/3\mathbb{Z} = Z(\mathrm{SU}(3))$ (the centre of
$\mathrm{SU}(3)$), with the three pairs as the fundamental
representation.
The full $\mathrm{SU}(3)$ acts as unitary rotations on the
three-dimensional pair space; the Casimir spectrum $(3, 5, 6)$
breaks $\mathrm{SU}(3) \to \mathrm{U}(1)^2$ (the Cartan subgroup).

\emph{Gravity $=$ colour}: on the base $\mathbf{H}^2$, the theory
is gravity (the graviton $j = 2$).
On the Galois cover
$(\mathbf{H}^2)^3/\mathbb{Z}_3$, the same theory has
$\mathrm{SU}(3)$ gauge symmetry: three copies of $\mathbf{H}^2$,
one per real embedding of~$K_7$, permuted by the Galois group.
The Galois correspondence IS the gauge--gravity duality.

\subsection{$\mathrm{SU}(2)$ from the Chern--Simons adjoint at $N = 4$}

\begin{proposition}[$\mathrm{SU}(2)$ from the adjoint Wilson line]
\label{prop:su2}
At $N = 4$, $m^* = 2$: $f(2,4) = 2 = j(j{+}1)$ with $j = 1$.
The $j = 1$ representation of $\mathrm{SL}(2,\mathbb{R})$ is the
\emph{adjoint} (dimension~$3$).
Wick rotation to the compact form gives $\mathrm{SU}(2)$; the
adjoint Wilson line carries $\mathrm{SU}(2)$ gauge symmetry.
\end{proposition}

\subsection{$\mathrm{U}(1)$ from Kaluza--Klein}

The $S^1$ fiber with flux $N/2$ gives a $\mathrm{U}(1)$ gauge field
at all~$N$, with coupling $\alpha = 6/(\pi c)$ and the exact lock
$\alpha/(8\pi G) = 1/(2\pi^2)$.

%% ====================================================================
\section{The Weinberg angle}
\label{sec:weinberg}

\begin{proposition}[Weinberg angle from the polygon hierarchy]
\label{prop:weinberg}
At $N = 4$, the critical mode $m^* = 2$ carries KK hypercharge
$Y = m^*/N = 1/2$ and normalised SU$(2)$ Casimir
$C_2/(N{-}1) = f(2,4)/3 = 2/3$.
The Weinberg angle is
\begin{equation}\label{eq:weinberg}
  \sin^2\theta_W = \frac{Y^2}{Y^2 + C_2/(N{-}1)}
  = \frac{1/4}{1/4 + 2/3} = \frac{3}{11} \approx 0.2727.
\end{equation}
\end{proposition}

\noindent
The experimental value at $M_Z$: $\sin^2\theta_W = 0.2312$.
The SU$(5)$ GUT prediction: $3/8 = 0.375$ (at the GUT scale).
Our prediction $3/11 = 0.273$ is 18\% above experiment,
versus SU$(5)$'s 62\%---substantially closer, and the remaining
discrepancy is expected from renormalisation-group running
from the polygon scale to~$M_Z$.

%% ====================================================================
\section{The two-force unification}
\label{sec:two-forces}

The four forces of the Standard Model reduce to two:

\paragraph{Force~1: Gravity/Colour ($N = 7$ sector).}
On the base $\mathbf{H}^2$: gravity (spin-$2$ graviton,
coupling~$G = 3/(2c)$).
On the Galois cover $(\mathbf{H}^2)^3 / \mathbb{Z}_3$:
$\mathrm{SU}(3)$ colour ($g_s^2 = 3 g_{\mathrm{CS}}^2$,
three colours from three Galois conjugates).
Long range: gravity (curvature).
Short range: colour (confinement $=$ compactification of the cover).
The Galois correspondence $K_7/\mathbb{Q}$ IS the
gauge--gravity duality for this sector.

\paragraph{Force~2: Electroweak ($N = 4$ sector $+$ KK).}
$\mathrm{SU}(2)$: from the CS adjoint Wilson line at $j = 1$ ($N = 4$).
$\mathrm{U}(1)$: from the $S^1$ KK photon.
Unified by the Seifert structure; the Weinberg angle
$\sin^2\theta_W = 3/11$ is determined by the KK hypercharge
and the CS Casimir.

\paragraph{Unification field.}
Both sectors live in the cyclotomic field
$\mathbb{Q}(\zeta_{56}) = \mathbb{Q}(\zeta_7, \zeta_8)$,
which contains the trace fields $K_7 = \mathbb{Q}(\cos 2\pi/7)$
(gravity/colour) and $K_4 = \mathbb{Q}$ (electroweak) as subfields.

%% ====================================================================
\section{UV completion}
\label{sec:uv}

The theory is UV-complete in the following precise sense.

\paragraph{Non-perturbative definition.}
The orbifold CFT at $c = N^2$---a $\mathbb{Z}_N$ orbifold of $N^2$
free bosons on $S^1$---is an exactly solvable 2D conformal field theory
with known partition function, OPE coefficients, and conformal blocks.
In the AdS$_3$/CFT$_2$ framework, this boundary CFT
\emph{defines} the bulk Chern--Simons theory at $k = N^2/12$
non-perturbatively.

\paragraph{Super-renormalisability.}
After KK reduction to $1{+}1$D, the superficial degree of divergence
is $D = 2 - E$ ($E$ = external legs).
Only the $2$-point function ($D = 0$) has a logarithmic divergence;
all vertices with $E \ge 4$ are UV-finite.
Two renormalisation parameters suffice: the vortex core size~$a$
(UV cutoff for self-energy) and the central charge~$c$ (overall
energy scale).

\paragraph{Physical observables.}
All physical observables---Havelock eigenvalues, partition function,
couplings, cosmological constant, Weinberg angle---are
\emph{exact} (no perturbative approximation).
The only divergences (self-energy at $d = 0$, KK zero-point sum)
are renormalised by physical parameters within the framework.

%% ====================================================================
\section{The fermion mass structure}
\label{sec:fermions}

\subsection{The Yukawa texture from KK charge conservation}

The Yukawa coupling $Y_{ij}$ between fermion pairs $i$, $j$ and the
Higgs (pair~3, modes $3$ and~$4$) is nonzero iff
$m_i - m_j + m_H \equiv 0 \pmod{7}$.
Evaluating for all pairs:
\begin{equation}\label{eq:yukawa-texture}
  Y = \begin{pmatrix}
    0 & Y_{12} & Y_{13} \\
    Y_{21} & Y_{22} & 0 \\
    Y_{31} & 0 & 0
  \end{pmatrix},
\end{equation}
with four texture zeros at $(1,1)$, $(2,3)$, $(3,2)$, $(3,3)$---determined
entirely by the arithmetic of $\mathbb{Z}/7\mathbb{Z}$.

\subsection{The 2+1 mass pattern and the exact tree-level ratio}

The matrix $YY^\dagger$ is block-diagonal:
a $2 \times 2$ block (pairs~1,2) and a $1 \times 1$ block (pair~3).
The eigenvalues of the $2 \times 2$ block
(in the large-$\sigma$ limit where the RS profiles saturate) are
$\lambda_\pm = (5 \pm \sqrt{13})/2$,
giving the exact tree-level mass ratio
\begin{equation}\label{eq:tree-ratio}
  \frac{m_1}{m_2}\bigg|_{\mathrm{tree}}
  = \sqrt{\frac{5 + \sqrt{13}}{5 - \sqrt{13}}}
  \approx 2.48,
\end{equation}
while the lightest mass $m_3 \propto e^{-\sigma/2}$ is
exponentially suppressed by the warp factor~$\sigma$.

\subsection{Up--down splitting from the $N = 4$ isospin}

The $\mathrm{SU}(2)$ doublet structure from the $N = 4$ sector
gives different KK masses for up-type ($\mu_4 = 1/2$) and
down-type ($\mu_4 = 3/2$) fermions.
The combined bulk mass $c_{k,I} = \sqrt{\mu_{7,k}^2 + \mu_{4,I}^2}$
separates the two types, with down-type having larger~$c$
(more UV-localised, hence lighter).

\subsection{The physical hierarchy from Yukawa RG running}

The tree-level ratio~$m_1/m_2 \approx 2.5$ is amplified
by the standard Yukawa $\beta$-function
($dy_t/d\ln\mu \propto y_t^3$) during running from the
polygon scale~$M_{\mathrm{poly}}$ to~$M_Z$.
An amplification factor of $\sim\!100$ produces
$m_t/m_c \approx 250$---close to the experimental~$280$.

\subsection{Sector decoupling}

The $N = 7$ and $N = 8$ polygon sectors are decoupled by the
mod-$56$ arithmetic of the common overfield
$\mathbb{Q}(\zeta_{56})$: the cross-sector Yukawa coupling
requires $8(m_1 + m_2) \equiv 28 \pmod{56}$, which has no
integer solution ($\gcd(8,56) = 8$ does not divide~$28$).
The texture zeros of the $N = 7$ Yukawa are topologically
protected.

%% ====================================================================
\section{The cosmological energy budget}
\label{sec:cosmology}

The free energy of the Havelock Field Theory at polygon number~$N$
decomposes into three components:
\begin{align}
  F_{\mathrm{DE}} &= b(N) + \tfrac{N}{4}
    & &\text{(vacuum $+$ radion zero-point)}, \label{eq:F-DE}\\
  F_{\mathrm{DM}} &= \tfrac{1}{2}\!\sum_{m \ne m^*}\!
    \ln|\lambda_m|
    & &\text{(frozen mode one-loop)}, \label{eq:F-DM}\\
  F_{\mathrm{M}}  &= \Delta E
    & &\text{(WDW mass gap = baryonic matter)}. \label{eq:F-M}
\end{align}
The radion contribution $N/4$ in~\eqref{eq:F-DE} is the zero-point
energy of the $S^1$ fiber modulus, stabilised at unit radius by
the flux~$N/2$: the radion mass is $M_{\mathrm{rad}} = N/2$,
giving $(1/2)M_{\mathrm{rad}} = N/4$.
This is \emph{not} a free parameter---it is determined by the geometry.

\begin{proposition}[Cosmological energy budget at $N = 11$]
\label{prop:cosmology}
At $N = 11$, $\Delta E \approx 0.81$:
\begin{center}
\small
\begin{tabular}{lccc}
\toprule
Component & Prediction & Planck 2018 & Tension \\
\midrule
Dark energy $\Omega_\Lambda$ & $68.9\%$ & $68.5 \pm 0.7\%$ & $0.6\sigma$ \\
Dark matter $\Omega_{\mathrm{DM}}$ & $26.9\%$ & $26.6 \pm 0.7\%$ & $0.4\sigma$ \\
$\Omega_\Lambda / \Omega_{\mathrm{DM}}$ & $2.56$ & $2.57$ & $0.5\%$ \\
Baryonic matter $\Omega_b$ & $4.2\%$ & $4.9 \pm 0.06\%$ & $12\sigma$ \\
\bottomrule
\end{tabular}
\end{center}
\noindent
The dark energy and dark matter fractions are within $1\sigma$ of the
Planck measurements.
The $\Omega_\Lambda/\Omega_{\mathrm{DM}}$ ratio matches to~$0.5\%$.
The baryon fraction has residual tension that depends
on the mass gap~$\Delta E$.
\end{proposition}

\subsection{Why $N = 11$}

$N = 11$ is distinguished in the polygon hierarchy:
\begin{itemize}
\item The trace field $K_{11} = \mathbb{Q}(\cos 2\pi/11)$ has
  degree~$[K_{11}:\mathbb{Q}] = 5 = \operatorname{rank}\mathrm{SU}(5)$,
  the grand unified group.
\item $N = 11$ sits at the \emph{quantum regime boundary}:
  the quantum phase diagram (\S\ref{sec:uv}) has classical
  $N \le 7$, marginal $8$--$10$, quantum $N \ge 11$.
\item There are $5$ palindromic mode pairs---the fundamental
  representation of $\mathrm{SU}(5)$.
\end{itemize}

\subsection{The dark matter spectrum}

The frozen Havelock modes are the dark matter candidates:
\emph{massive} (from the Casimir gaps), \emph{stable}
($\mathbb{Z}_N$-protected), \emph{non-luminous} (different
KK charge from the Higgs), and \emph{cold} (zero kinetic energy
at threshold).
At $N = 11$, the Casimir gaps are $\lambda_m = 1, 3, 6, 10$
(the \emph{triangular numbers} $T_1, T_2, T_3, T_4$),
giving four dark matter species with mass ratios
\begin{equation}\label{eq:dm-masses}
  M_1 : M_2 : M_3 : M_4
  = 1 : \sqrt{3} : \sqrt{6} : \sqrt{10}.
\end{equation}
There are $8$ dark matter particles in total
($4$ frozen pairs $\times$ degeneracy~$2$).

\subsection{The baryon fraction and the mass gap}

The baryon fraction within total matter:
$\Omega_b / (\Omega_b + \Omega_{\mathrm{DM}})
= F_M / (F_{\mathrm{DM}} + F_M)$.
At $N = 11$: $F_{\mathrm{DM}} = 5.193$.
The Planck value $\Omega_b/\Omega_m = 15.6\%$ requires
$F_M = 0.96$---slightly above the WDW estimate
$\Delta E \approx 0.81$.
\subsection{The mass gap: $\Delta E \to \sqrt{2/\pi}$}

The WDW mass gap, computed via non-uniform FEM mesh
(quadratically concentrated near $\rho = 0$), gives converged values:
\begin{center}
\small
\begin{tabular}{lcccccc}
\toprule
$N$ & $6$ & $7$ & $8$ & $9$ & $10$ & $11$ \\
\midrule
$\Delta E$ & $0.841$ & $0.832$ & $0.825$ & $0.821$ & $0.817$ & $0.815$ \\
\bottomrule
\end{tabular}
\end{center}
\noindent
The asymptotic formula (fit to $< 0.05\%$ error):
\begin{equation}\label{eq:mass-gap}
  \Delta E = \sqrt{\frac{2}{\pi}}
  + \frac{0.45\,\ln c}{c} + O(1/c),
\end{equation}
where $\sqrt{2/\pi} = 0.79788\ldots$ is the mean of the
half-normal distribution.
At $N = 11$: $\Delta E = 0.815$, giving $\Omega_b = 4.2\%$
(observed: $4.9\%$, a $14\%$ discrepancy).
The remaining tension requires either finite-temperature
corrections or higher-order loop contributions to~$F_{\mathrm{DM}}$.

%% ====================================================================
\section{Discussion}
\label{sec:discussion}

The Havelock Field Theory is specified by a single integer~$N$
on a single geometry $\mathbb{R} \times (\mathbf{H}^2 \times_N S^1)$.
It produces:
\begin{itemize}
\item Gravity $+$ $\mathrm{SU}(3) \times \mathrm{SU}(2)
  \times \mathrm{U}(1)$ gauge symmetry;
\item $\Lambda = (N^2 - 16)/16$ ($\Lambda = 0$ at $N = 4$,
  positive for $N \ge 5$);
\item $\sin^2\theta_W = 3/11$ (18\% from experiment);
\item Exact coupling lock $\alpha/(8\pi G) = 1/(2\pi^2)$;
\item Zero Weyl anomaly on Einstein backgrounds;
\item UV completion via orbifold CFT;
\item The graviton at $N = 7$ (unique);
\item Two forces: gravity/colour $+$ electroweak.
\end{itemize}

\noindent
What it does \emph{not} produce:
the fermion mass hierarchy, the CKM matrix,
the number of generations from first principles
(though $\mathrm{ind}(D) = N/2 + (1{-}g) = 3$ for
$(N,g) = (4,0)$, $(6,1)$, or $(8,2)$),
or a resolution of the landscape ($N$ is an input, not derived).

\medskip\noindent
The theory is built from four classical ingredients:
Havelock (1931), Kaluza--Klein (1921),
Thurston (1982), and Chern--Simons (1989).
No string theory. No extra dimensions beyond~$S^1$.

\bibliographystyle{plainnat}
\bibliography{../shared/refs}
\end{document}
